\documentclass[preprint,11pt,authoryear]{elsarticle}

\usepackage{amsmath,amssymb,amsthm}
\usepackage{booktabs}
\usepackage{tikz}
\usepackage{pgfplots}
\pgfplotsset{compat=1.18}
\usepackage[margin=1in]{geometry}
\usepackage[hidelinks]{hyperref}
\hypersetup{pdftitle={Credit Lines for Pseudonymous Borrowers: Identity-Gated Bayesian Lending under an On-Chain Issuance Budget}, pdfauthor={Claude Code}}

\newtheorem{proposition}{Proposition}
\newtheorem{lemma}{Lemma}

\makeatletter
\newcommand{\inputrows}[1]{\@@input #1 }
\makeatother

% Generated by scripts/extract_issuer_policy.py; do not edit.
\newcommand{\CalCycles}{53{,}935}
\newcommand{\CalDefaults}{217}
\newcommand{\CalPredicted}{0.43}
\newcommand{\CalTrue}{0.41}
\newcommand{\CalRealised}{0.40}
\newcommand{\CalLo}{0.35}
\newcommand{\CalHi}{0.46}

% Generated by scripts/extract_issuer_policy.py; do not edit.
\newcommand{\LineShiftBasic}{-0.29}
\newcommand{\LineShiftRelBasic}{-4.8}
\newcommand{\LineShiftStandard}{-0.21}
\newcommand{\LineShiftRelStandard}{-1.0}
\newcommand{\LineShiftEnhanced}{-0.14}
\newcommand{\LineShiftRelEnhanced}{-0.3}
\newcommand{\LineShiftMaxAbs}{0.29}
\newcommand{\LineShiftMaxRel}{4.8}
\newcommand{\PremiumPerCycle}{0.658}
\newcommand{\DemandMean}{50}

% Generated by scripts/misspecification.py; do not edit.
\newcommand{\BrierWeightedAssumed}{0.447}
\newcommand{\SkillWeightedAssumed}{0.0}
\newcommand{\RatioWeightedAssumed}{1.05}
\newcommand{\FarmLineWeighted}{6.10}
\newcommand{\MeanLineWeightedAssumed}{23.6}
\newcommand{\BrierUnweightedAssumed}{0.446}
\newcommand{\SkillUnweightedAssumed}{0.0}
\newcommand{\RatioUnweightedAssumed}{1.10}
\newcommand{\FarmLineUnweighted}{15.00}
\newcommand{\MeanLineUnweightedAssumed}{26.5}
\newcommand{\BrierTierOnlyAssumed}{0.478}
\newcommand{\SkillTierOnlyAssumed}{0.0}
\newcommand{\RatioTierOnlyAssumed}{1.05}
\newcommand{\FarmLineTierOnly}{6.10}
\newcommand{\MeanLineTierOnlyAssumed}{19.7}
\newcommand{\RatioWeightedRiskier}{1.42}
\newcommand{\RatioLateWeightedRiskier}{1.15}
\newcommand{\SkillWeightedRiskier}{-0.0}
\newcommand{\RatioTierOnlyRiskier}{1.45}
\newcommand{\RatioWeightedLateHalf}{1.09}
\newcommand{\RatioLateWeightedLateHalf}{0.81}
\newcommand{\SkillWeightedLateHalf}{0.0}
\newcommand{\RatioTierOnlyLateHalf}{1.05}
\newcommand{\RatioWeightedLateDouble}{1.00}
\newcommand{\RatioLateWeightedLateDouble}{0.77}
\newcommand{\SkillWeightedLateDouble}{0.0}
\newcommand{\RatioTierOnlyLateDouble}{1.05}
\newcommand{\RatioWeightedShift}{1.50}
\newcommand{\RatioLateWeightedShift}{1.74}
\newcommand{\SkillWeightedShift}{-0.0}
\newcommand{\RatioTierOnlyShift}{1.53}
\newcommand{\RatioWeightedMin}{1.00}
\newcommand{\RatioWeightedMax}{1.50}
\newcommand{\MisspecSeeds}{3}
\newcommand{\MisspecHonest}{600}


\journal{European Journal of Operational Research}

\begin{document}

\begin{frontmatter}

\title{Credit Lines for Pseudonymous Borrowers: Identity-Gated Bayesian Lending under an On-Chain Issuance Budget}

\author{Claude Code\fnref{fn1}}
\fntext[fn1]{Claude Code is an AI agent made by Anthropic, working with the Microcredit Protocol
Project; the research was directed by the project's human operator. Sources, data and code:
\url{https://github.com/scottonchain/microcredit-theory}.}
\affiliation{organization={Microcredit Protocol Project}, country={International}}

\begin{abstract}
A lending pool on a public blockchain lends without collateral against credit lines that an issuer
publishes for pseudonymous accounts. The pool's contract bounds what the issuer can put at risk: a
budget on the sum of lines, charged on each account's highest line since the line was last unused,
and a cap on how much one report may add. The issuer must then decide whom to lend to and how much,
from thin and manipulable evidence. We give a policy that does so in four steps. Only accounts with
a costly identity are eligible, and a tier's cap never exceeds the cost of a fraudulent identity in
it, which bounds the profit from buying identities by their cost, outside two channels the analysis
names. Repayment history enters a Beta-Binomial estimate of the probability of default only in
proportion to the risk lenders actually bore, so every route to evidence that costs nothing earns
nothing. Each line maximises a first-order surrogate of the pool's expected surplus under exponential
demand, at $\mu\ln(i/p)$. Lines are allocated within the budget and the per-report cap by a greedy
rule that is optimal because the constraints form a laminar matroid, and every report is planned to
land even if an adversary releases held budget first. In a 24-month simulation the contract accepted
all 925 reports, predicted default rates matched realised ones in every decile of the world the
prior was fitted to, lines grew for low-risk borrowers and shrank for high-risk ones, and attackers
buying identities lost money in every tier unless identities cost far less than the issuer assumed.
Scored on a common panel independent of the lending decision, the repayment evidence improves on
the tier prior at this horizon by a measurable but small amount, and the weights do not improve on
an unweighted rule; their contribution is to close the routes to evidence that cost nothing, which
an unweighted rule leaves open.
\end{abstract}

\begin{keyword}
OR in banking \sep credit scoring \sep Bayesian estimation \sep resource allocation \sep matroids \sep incentive compatibility
\end{keyword}

\end{frontmatter}

% ---------------------------------------------------------------------------------------------
\section{Introduction}

Credit scoring decides who receives credit and how much \citep{thomas2017credit}. The operational
research literature has studied it as classification \citep{lessmann2015benchmarking} and as a
decision problem in which the lender sets limits to maximise profit
\citep{crook2007recent}. This paper studies a setting that changes three of the standing
assumptions. Borrowers are pseudonymous accounts on a public blockchain, so an attacker can create as
many as it likes \citep{douceur2002sybil}. Their history is public and partly manufacturable: a loan
repaid within a day pays no interest, and capital returned at the end of a history can be reused
for the next account. And the lender, here called the issuer, does not lend itself. It publishes
lines into a pool funded by others, whose contract limits how much the issuer can put at risk.

The pool's contract charges an issuance budget on each account's highest line since the line was
last unused, releases that budget only when the account has no open loans and no commitments, and
caps how much one report may raise the total \citep{conservedcredit2026}. Under those rules, the
lenders' potential loss on the issuer's lines never exceeds the budget, even if the issuer is
compromised. Within that bound, the issuer still has to be accurate and resistant to manipulation:
it is a delegated monitor in the sense of \citet{diamond1984delegated}, and its lines are where the
pool's losses come from.

We give a reference policy for this issuer and analyse its properties. The contributions are:
\begin{itemize}
\item \emph{An incentive-compatible eligibility rule} (Proposition~\ref{prop:ic}). A tier's line
  cap must not exceed the cost of an identity in that tier. Then no number of bought identities,
  with any history, yields a profit, because the loss an account can cause over its lifetime is at
  most its highest line plus the dues it paid (Lemma~\ref{lem:lifetime}).
\item \emph{Evidence weighted by risk borne} (Section~\ref{sec:evidence}). A repayment counts in a
  Beta-Binomial posterior only in proportion to the unsecured principal, the time at risk and its
  recency. Every zero-cost route to evidence has weight zero.
\item \emph{A closed-form line} (Proposition~\ref{prop:line}): with exponential demand, expected
  profit is concave in the line and maximised at $\mu\ln(i/p)$.
\item \emph{Optimal allocation within the contract's rules} (Proposition~\ref{prop:greedy} and
  Lemma~\ref{lem:frontrun}). The feasible raises form a laminar matroid, so a greedy rule is optimal;
  and the plan respects every rule even when an adversary front-runs it.
\item \emph{A simulation} of 24 monthly reports with honest borrowers, identity buyers, a farm of
  recycled histories and an adversary who front-runs every report.
\end{itemize}

% ---------------------------------------------------------------------------------------------
\section{The Setting}\label{sec:setting}

\paragraph{The pool} Accounts borrow stablecoins for 30-day terms. An account's own credit is its
issued line plus the dues it has paid, a share of its interest held in a first-loss reserve; other
accounts may also back it with credit of their own \citep{conservedcredit2026}. A loan defaults 30
days after its due date and is then written off.

\paragraph{The issuer's interface} The issuer publishes scores $q(a)\in[0,1]$; the line is
$\ell(a)=M\,q(a)$ for a maximum line $M$. The contract keeps a held budget $\beta(a)$ per account.
A report that raises $q(a)$ above $\beta(a)$ raises $\beta(a)$ to it, and a report is rejected if
it raises $\sum_a\beta(a)$ above a budget $B$ or by more than a per-report cap $\Delta$. Anyone may
call a release that lowers $\beta(a)$ to $q(a)$, but only while $a$ has no open loans and no backing
commitments. Scores read as zero if no report arrives within a maximum age, so the issuer sends
empty reports as heartbeats. Reports carry increasing epochs and at most 500 accounts each.

\paragraph{What the issuer reads} For each account: whether an identity provider has verified it,
whether it has ever defaulted, its own loans (principal, term, disbursement, repayment and closing
times) and, per loan, the secured backing it held. From its own registry: the identity behind the
account and the identity's tier. It never reads who backs whom. A function of the backing graph is
exactly what Sybil accounts can manipulate \citep{cheng2005sybilproof}.

% ---------------------------------------------------------------------------------------------
\section{Eligibility and Incentive Compatibility}\label{sec:ic}

An account receives a line only if it holds a verified identity of tier $t$, obtained at a cost
$k_t$ (the price of a fraudulent identity in that tier, plus any expected penalty), uses the
identity's one live address, and has no default on any address of the identity. Everything else
receives zero. The lifetime constraint is on the identity, not the address: in the implementation an
identity is bound to one address for life and cannot move, so the address's record is the identity's
record. If a later version lets an identity rotate to a new address, the credit charged to the old
address as a backer and its default flag must carry over with its repayment history; one live address
at a time is not by itself the lifetime invariant that Lemma~\ref{lem:lifetime} needs.

\begin{lemma}[Lifetime loss per account]\label{lem:lifetime}
For any account $a$ and any sequence of operations, the losses lenders bear that are attributable to
$a$, namely the credit charged to $a$ as a backer plus the residual of its own defaults, are at most
$\hat\ell(a)+d(a)$, where $\hat\ell(a)$ is the highest line $a$ has ever held and $d(a)$ its dues.
\end{lemma}

\begin{proof}
The per-account invariant of \citet[Lemma~5.1]{conservedcredit2026} states that residual exposure
plus committed credit plus attributed loss is at most the line plus dues, and each operation of the
pool preserves it. With the line replaced by its running maximum $\hat\ell(a)$, every step of that
proof goes through unchanged: borrowing and committing are bounded by granted credit, which is at
most $\hat\ell(a)+d(a)$ less the credit already charged; a charge lowers commitments by at least
what it adds to the loss; a default moves the residual from exposure into the loss; and lowering a
line leaves $\hat\ell(a)$ unchanged. Since the left side of the invariant includes the attributed
loss, the claim follows.
\end{proof}

\begin{proposition}[Incentive compatibility]\label{prop:ic}
Let an attacker buy $n_t$ identities in each tier $t$ at cost $k_t$ and then follow any strategy,
possibly holding a share $s$ of the pool's shares throughout as a lender. Measure its
\emph{deviation profit} as its final liquid wealth less its initial wealth, less what the same
initial wealth would have earned as a passive lender holding the same share $s$ over the same
period; identities and accounts it abandons are worth nothing to it. Then its deviation profit is at
most
\[
  \sum_t n_t\bigl(\min\{\mathrm{cap}_t,M\}-k_t\bigr)\;+\;I\bigl(r-1+s(1-f-r)\bigr)\;+\;s\,r\,I\,\mathbb{1}_{\mathrm{stress}}\;+\;H,
\]
where $I$ is the interest its accounts paid, $r$ and $f$ the reserve and fee shares of interest,
$\mathbb{1}_{\mathrm{stress}}$ indicates that other borrowers' defaults consumed its dues from the
pooled reserve before it withdrew, and $H$ is the backing that honest accounts committed to it and
lost. If $\mathrm{cap}_t\le k_t$ in every tier, the first term is at most zero. The second is at most
$-fI\le0$. The third is at most $s\,r\,I$, and the second and third together are positive only if
$s>(1-r)/(1-f)$. The last is the honest backers' own loss, which no issuer rule controls.
\end{proposition}

\begin{proof}
By Lemma~\ref{lem:lifetime}, each identity yields at most its highest line, which the policy keeps
at or below $\min\{\mathrm{cap}_t,M\}$, plus its dues $d=rI$ in total. The attacker paid $k_t$ per
identity and interest $I$. Relative to the passive lender, it recovers through its shares at most
$s(1-f-r)I$ of that interest \citep[Proposition~7.4]{conservedcredit2026}, which gives the second
term; and if other borrowers' defaults consume its dues from the pooled reserve before it withdraws,
those dues absorb losses that would otherwise have lowered the value of its shares, which adds at
most $s\,rI$ \citep[Proposition~7.5]{conservedcredit2026}. Honest backing is the honest backers' own
risk, charged to their credit.
\end{proof}

The benchmark matters. Against total wealth, ordinary returns on passive lending would enter, and
the statement would be about the pool's yield rather than about identities; against passive lending,
the proposition says that buying identities, with any history, adds nothing to what the attacker
would have earned anyway, outside two channels it names. The first, the pooled reserve under stress,
is the channel of \citet[Proposition~7.5]{conservedcredit2026}: it pays only an agent that holds more
than $(1-r)/(1-f)$ of the pool, 70\% at the deployed reserve share of 30\% and no fee, and a run on a
copy of the deployed contract found an agent holding 91\% of the pool ending $0.49$ ahead of a
passive lender at $I=2.33$, as the formula predicts. The second, honest backing, is a loss that
honest backers choose to risk. So the proposition is not a claim that no deviation is ever
profitable; it is a bound on the profit from identities, and a statement of where any other profit
must come from. Under the condition, and outside those channels, a profit-seeking attacker gains
nothing from any number of identities through any history; this is the identity-cost form of the
result of \citet{friedman2001cheap} that, with cheap pseudonyms, newcomers must pay dues, and it is
why identities issued once per person matter \citep{borge2017personhood}. An attacker that does not
seek profit can still cost lenders up to $\sum_tn_t\,\mathrm{cap}_t$, which under the condition is at
most what it spent on identities, and at any moment at most $MB$ on all of the issuer's lines.
Table~\ref{tab:tiers} lists the tiers used below; their costs are illustrative inputs, which an issuer
must set from the market price of identities and revisit.

\begin{table}[ht]
  \centering\small
  \caption{Tiers: identity cost, line cap, the prior's strength in trials and its annual PD, a new
  identity's line, and the line after 24 clean 30-day trials (USDC).}
  \label{tab:tiers}
  \begin{tabular}{lrrrrrr}
    \toprule
    Tier & Cost $k_t$ & Cap & Prior strength & Prior PD & Starting line & After 24 trials \\
    \midrule
    \inputrows{data/table_tiers.tex}
    \bottomrule
  \end{tabular}
\end{table}

% ---------------------------------------------------------------------------------------------
\section{Evidence and the Probability of Default}\label{sec:evidence}

Each closed loan of the account itself is a fractional Bernoulli trial. A loan repaid on time is a
good trial with weight
\[
  w_{\mathrm{good}}=2^{-\mathrm{age}/h}\cdot\min\Bigl\{1,\frac{\text{principal}-\text{secured}}{P_{\mathrm{ref}}}\Bigr\}\cdot\min\Bigl\{1,\frac{\min\{\text{elapsed},\text{term}\}}{30\ \text{days}}\Bigr\},
\]
the product of recency with half-life $h$, unsecured size up to a reference principal
$P_{\mathrm{ref}}$, and time at risk; and $w_{\mathrm{good}}=0$ if the loan closed within the
24-hour window in which no interest accrues. A loan repaid after its due date is a bad trial whose
weight is not discounted for stake, so stake cannot hide lateness. Open loans are not yet evidence,
and a default ends the identity. Loans of other accounts, including accounts this one backed, are
never read.

Let $G$ and $B$ be the sums of good and bad weights. With a Beta$(\alpha_t,\beta_t)$ prior on the
rate $\theta$ of bad cycles, fitted by moments to the tier's population, the posterior is
Beta$(\alpha_t+B,\beta_t+G)$ \citep{gelman2013bda}. Down-weighting observations by a power of their
likelihood is the standard way to discount less informative data \citep{ibrahim2000power}. A bad
cycle is a default with probability $1/(1+c)$, where $c$ is the issuer's ratio of late repayments to
defaults, so the probability of default per 30-day cycle is $p=\mathbb{E}[\theta]/(1+c)$.

Table~\ref{tab:routes} shows how each weight closes a route to evidence that would otherwise cost
nothing. Interest-free loans and fully secured loans put nothing at risk, and secured backing can be
recycled through account after account \citep[Theorem~7.2]{conservedcredit2026}. A guarantor record
can be manufactured by backing one's own fresh account. Dust loans and very short loans would give
many trials for almost no interest. With the weights, evidence costs interest at a fixed rate, about
one unit per full trial at the protocol's loan rate. None of the weights is needed for
Proposition~\ref{prop:ic}, which rests on the caps alone; they keep the estimate honest, so that lines
go to accounts lenders have seen repay at risk.

\begin{table}[ht]
  \centering\small
  \caption{Routes to evidence, their weights as good and bad trials, and the interest paid per unit
  of good evidence (USDC; ``none'' when the route yields no good evidence).}
  \label{tab:routes}
  \begin{tabular}{p{0.52\linewidth}rrr}
    \toprule
    Route & Good & Bad & Interest per good trial \\
    \midrule
    \inputrows{data/table_routes.tex}
    \bottomrule
  \end{tabular}
\end{table}

% ---------------------------------------------------------------------------------------------
\section{The Line}\label{sec:line}

In normal times a borrower draws $\min\{\ell,D\}$ of its line, with $D$ exponential of mean $\mu$.
Before a default it draws the whole line, with loss given default of 100\%. With $i$ the risk premium
per 30-day cycle, expected issuer profit per cycle is premium on the expected draw less expected loss.

\begin{proposition}[Profit-maximising line]\label{prop:line}
$\Pi(\ell)=i\,\mu\bigl(1-e^{-\ell/\mu}\bigr)-p\,\ell$ is strictly concave, and it is maximised at
\[
  \ell^*=\mu\ln(i/p)\ \text{ if } p<i,\qquad \ell^*=0\ \text{ otherwise}.
\]
\end{proposition}

\begin{proof}
$\mathbb{E}\min\{\ell,D\}=\int_0^\ell e^{-x/\mu}\,dx=\mu(1-e^{-\ell/\mu})$. The marginal profit
$\Pi'(\ell)=i\,e^{-\ell/\mu}-p$ is strictly decreasing, positive at zero exactly when $p<i$, and zero
at $\mu\ln(i/p)$.
\end{proof}

The marginal unit of line earns the premium only when it is drawn, but adds to the loss at a default,
when it is always drawn; so lines are rationed below demand, as in models of credit rationing
\citep{stiglitz1981rationing}. The policy's target is $\min\{\ell^*,\mathrm{cap}_t\}$, and an account
with an overdue loan gets zero at once.

\paragraph{What the objective is} $\Pi$ is a first-order surrogate, and its economic interpretation
should be stated. Normal demand is specified conditional on not defaulting, and a borrower that
defaults draws the whole line and pays no premium that cycle, so the cash-flow objective per cycle
on a line, premium income less principal loss, is
\[
  \tilde\Pi(\ell)=(1-p)\,i\,\mu\bigl(1-e^{-\ell/\mu}\bigr)-p\,\ell,
\]
strictly concave with optimum $\mu\ln\bigl((1-p)\,i/p\bigr)$, lower than $\ell^*$ by $\mu\ln(1-p)$.
Table~\ref{tab:objective} gives the difference at each tier's prior: at most \LineShiftMaxAbs\ USDC,
or \LineShiftMaxRel\% of the line, in the basic tier, where the line is smallest, and within the
one-unit allocation step in the other two. The funding rate on a defaulted balance over the late
period, fees, utilisation and the reserve's share add further terms, which separate the pool's gross
loan surplus from lenders' net return. Whose objective it is also needs saying: the issuer does not
lend. $\Pi$ is the pool's expected gross surplus per cycle on the line, which lenders and the
first-loss reserve share in the proportions $1-r$ and $r$ of interest; the issuer maximises it as the
pool's agent. It is not lenders' net return, and the companion pricing paper
\citep{pricingreserve2026} shows lenders earning less than the funding rate at the reserve share of
65\% used in the simulation below, because the reserve withholds part of the surplus. The two
statements coexist: a positive surplus does not ensure that lenders earn their alternative return,
and the choice of reserve share is the pricing paper's subject, not this one's.

\begin{table}[ht]
  \centering\small
  \caption{The surrogate objective against the cash-flow objective at each tier's prior PD per
  cycle ($\mu=\DemandMean$, $i=\PremiumPerCycle\%$ per cycle): the optimal line under each (USDC)
  and the difference.}
  \label{tab:objective}
  \begin{tabular}{lrrrrr}
    \toprule
    Tier & $p$ per cycle & $\mu\ln(i/p)$ & $\mu\ln((1-p)i/p)$ & difference & relative \\
    \midrule
    \inputrows{data/table_objective.tex}
    \bottomrule
  \end{tabular}
\end{table}

% ---------------------------------------------------------------------------------------------
\section{Allocation within the Budget}\label{sec:allocation}

Lowering a line is free and immediate. Raises are split into segments of one unit of line, and the
segments of account $a$ between its current score $S_a$ and its held budget $H_a$ use only the
per-report cap, while those above $H_a$ also use new budget. Segment $j$ of account $a$ has marginal
profit $\Pi_a'(x_j)$, which decreases in $j$ by Proposition~\ref{prop:line}.

\begin{proposition}[Greedy allocation]\label{prop:greedy}
Taking segments in decreasing order of marginal profit, skipping any that no longer fits the
per-report cap or the remaining budget, and stopping at non-positive marginal profit, maximises total
expected profit over all feasible sets of raises.
\end{proposition}

\begin{proof}
Every segment uses one unit of the per-report cap, and the segments above held budget also use one
unit of budget. The feasible sets are those with at most $\Delta$ segments in all and at most the
remaining budget among the segments above held budget: constraints on a laminar family of sets, so
they are the independent sets of a laminar matroid. The greedy algorithm finds a maximum-weight
independent set of any matroid \citep{edmonds1971matroids}. Within an account, segments below held
budget have larger marginal profit than those above it, and marginal profit decreases along each
account's segments, so the greedy order takes each account's segments as a prefix, which is what a
raise to a single new score requires.
\end{proof}

The problem is a two-constraint fractional knapsack \citep{dantzig1957discrete}. On 200 random
instances with both constraints binding, the greedy matched the linear-programming optimum computed
with HiGHS \citep{huangfu2018highs} within $10^{-6}$ relative error.

A plan is computed from the state the issuer reads, but the release can change that state before the
report lands, and anyone may call it.

\begin{lemma}[Front-running]\label{lem:frontrun}
Between the issuer's read and its report, only a release can change the provider's state, and it only
lowers held budget from $H_a$ to $S_a$. For an account in the report with new score $N_a$, a release
followed by the report changes the held total by $-(H_a-S_a)+\max\{0,N_a-S_a\}\le\max\{0,N_a-H_a\}$,
and raises the report's increase for that account to at most $\max\{0,N_a-S_a\}$.
\end{lemma}

\begin{proof}
If $N_a\le H_a$ the release lowers the held total by $H_a-S_a$ and the report raises it by at most
$\max\{0,N_a-S_a\}\le H_a-S_a$. If $N_a>H_a$ the change is $N_a-H_a$. The increase counted against the
per-report cap is measured from the held budget at the time the report lands, which is at least
$S_a$.
\end{proof}

The policy therefore plans every report so that it lands in both states, as read and with every
held budget released: $\sum_a\max\{0,N_a-S_a\}\le\Delta$ per cycle across all batches;
$\sum_a\beta(a)+\sum_a\max\{0,N_a-H_a\}\le B$; no raise at all when governance has lowered $B$
below the held total; scores integral in $[0,\min\{\mathrm{cap}_t,M\}]$; at most 500 accounts per
report; and consecutive epochs. A dry run against both copies of the provider's state stops the cycle
if either would revert.

% ---------------------------------------------------------------------------------------------
\section{Simulation}\label{sec:sim}

\paragraph{World} 24 monthly reports, maximum line 100, loan rate 12.33\% (a funding rate of 4.33\%
and a premium of 800 basis points), reserve share 65\%: the production deployment script's defaults
at the pinned contract commit, where the share has since become an interim 45\% (6 October 2026);
the live testnet deployment uses 500 basis points and 30\%. Honest verified borrowers are 40\% basic,
40\% standard and 20\% enhanced, with true annual PDs of 1, 2, 5, 10 or 20\% drawn from their tier's
population, the population the prior is fitted to. Each cycle they borrow with probability 0.85,
default with the cycle PD after drawing their whole limit, repay late with twice that probability,
and otherwise repay an exponential draw of mean 50 on time; a quarter of on-time repayments land
between the issuer's read and its report. A griefer releases every held budget just before each
report lands. Identity buyers repay a reference loan on time each month and draw everything at their
best month. A farm runs the recycled-seed construction on unverified accounts and on one bought
identity. Draws are seeded per borrower, so runs that differ only in budget or attackers face the
same borrowers.

\paragraph{The contract's rules} All 925 reports planned in the base run and the budget sweep were
accepted by an exact port of the provider's contract, each after a front-run release (54 lines
released in total). The per-report cap and the budget each reached full use in some cycle, and
unsecured exposure on the issuer's lines peaked at 83\% of held budget plus dues.

\paragraph{Farms} An unverified account that repaid 96 loans and 9,600 units inside the
interest-free window received a line of 0; a rule granting 25\% of repaid principal would have given
it 100. The same farm on a bought basic identity received 6.10, exactly the line of a new basic
identity with no history.

\paragraph{Honest lines} Lines start small and grow with repayment, the progressive lending of
microfinance \citep{morduch1999promise,armendariz2010}. Month 0 is small because the per-report cap
rations the first report. After the ramp, lines rise for low-risk borrowers and fall for high-risk
ones as lateness shows (Table~\ref{tab:lines}, Figure~\ref{fig:lines}).

\begin{table}[ht]
  \centering\small
  \caption{Mean line of surviving honest borrowers by true annual PD (USDC; base run of 1,000
  borrowers, budget of 250 lines, 60 per report) and the share alive at month 23.}
  \label{tab:lines}
  \begin{tabular}{rrrrrr}
    \toprule
    True PD & Month 0 & Month 6 & Month 12 & Month 23 & Alive at month 23 \\
    \midrule
    \inputrows{data/table_lines.tex}
    \bottomrule
  \end{tabular}
\end{table}

\begin{figure}[ht]
  \centering
  % Figure: mean line of surviving honest borrowers by month (data/honest_lines.csv).
\begin{tikzpicture}
  \begin{axis}[
      width=0.72\linewidth, height=5.4cm,
      xmin=0, xmax=23, ymin=0, ymax=32,
      xlabel={month}, ylabel={mean line (USDC)},
      tick label style={font=\footnotesize}, label style={font=\small},
      legend style={font=\scriptsize, draw=none, at={(0.5,1.03)}, anchor=south},
      legend columns=5, legend cell align=left,
    ]
    \addplot[thick, green!45!black] table[x=month, y=pd1, col sep=comma] {data/honest_lines.csv};
    \addlegendentry{true PD 1\%}
    \addplot[thick, blue!65!black] table[x=month, y=pd2, col sep=comma] {data/honest_lines.csv};
    \addlegendentry{2\%}
    \addplot[thick, orange!80!black] table[x=month, y=pd5, col sep=comma] {data/honest_lines.csv};
    \addlegendentry{5\%}
    \addplot[thick, violet] table[x=month, y=pd10, col sep=comma] {data/honest_lines.csv};
    \addlegendentry{10\%}
    \addplot[thick, red!70!black] table[x=month, y=pd20, col sep=comma] {data/honest_lines.csv};
    \addlegendentry{20\%}
  \end{axis}
\end{tikzpicture}

  \caption{Mean line of surviving honest borrowers by month and true annual PD.}
  \label{fig:lines}
\end{figure}

\paragraph{Calibration} Over \CalCycles{} honest borrower-cycles in three seeds, the mean predicted
default probability per cycle is \CalPredicted\%, the mean true probability \CalTrue\%, and the
realised rate \CalRealised\% (\CalDefaults{} defaults; 95\% interval \CalLo{} to \CalHi\%). In every
decile of the prediction, the mean true probability tracks the prediction within 0.05 percentage
points, and the prediction lies inside the realised rate's 95\% interval (Figure~\ref{fig:calibration}).

\begin{figure}[ht]
  \centering
  % Figure: calibration by decile (data/calibration.csv; percentages per 30-day cycle).
\begin{tikzpicture}
  \begin{axis}[
      width=0.6\linewidth, height=6cm,
      xmin=0.2, xmax=0.7, ymin=0, ymax=0.95,
      xlabel={mean predicted default probability per cycle (\%)},
      ylabel={realised default rate (\%)},
      tick label style={font=\footnotesize}, label style={font=\small},
      scaled ticks=false,
    ]
    \addplot[dashed, gray, domain=0.2:0.7] {x};
    \addplot[only marks, mark=*, mark size=1.6pt, blue!65!black,
             error bars/.cd, y dir=both, y explicit]
      table[x=predicted, y=realised, y error plus expr=\thisrow{hi}-\thisrow{realised},
            y error minus expr=\thisrow{realised}-\thisrow{lo}, col sep=comma] {data/calibration.csv};
  \end{axis}
\end{tikzpicture}

  \caption{Realised default rate per cycle with its 95\% interval, against the mean predicted
  probability, by decile of the prediction. The dashed line is perfect calibration.}
  \label{fig:calibration}
\end{figure}

\paragraph{Budget} Table~\ref{tab:budget} varies the budget with 600 honest borrowers. A tight
budget goes to the best marginal profit, so its loss rate is lowest; as it grows, riskier lines are
funded until demand rather than the budget binds, near 12,500 units. Loss per exposure-year exceeds
the line-weighted PD because a defaulter draws its whole line while others draw less. Realised profit
stays within about 20\% of the model's expectation.

\begin{table}[ht]
  \centering\small
  \caption{Issuer outcomes against the budget (USDC; 600 honest borrowers, mean of three seeds).}
  \label{tab:budget}
  \footnotesize\setlength{\tabcolsep}{4pt}
  \begin{tabular}{rrrrrrrr}
    \toprule
    & Mean & & Loss per & Loss per & Line-weighted & & Model \\
    Budget & held & Exposure & exposure-year & budget-year & PD & Profit & profit \\
    \midrule
    \inputrows{data/table_budget.tex}
    \bottomrule
  \end{tabular}
\end{table}

\paragraph{Validation under misspecification} The calibration above is measured in the world the
prior was fitted to, with lateness generated at the ratio the policy assumes; it shows that the
machinery is coherent there, not that the evidence weights predict better or resist a world that
differs. Tables~\ref{tab:misspec} and~\ref{tab:misspecpanel} run three scoring rules on the same
borrower paths in five worlds, \MisspecHonest\ honest borrowers and \MisspecSeeds\ seeds each: the
weights as proposed; an unweighted rule in which every repaid loan is one trial, good or bad,
whatever it put at risk; and the tier prior alone. The worlds are the assumed one; a population with
more of every tier at the high default probabilities than the prior assumes; lateness at half and at
twice the assumed ratio to defaults; and a regime shift in which every borrower's default and late
probabilities double from month 12. The scoring rules are compared by the Brier score of the
predicted cycle default probability against the realised default, by its skill against the realised
base rate, and by the ratio of realised to predicted defaults over all cycles, over months 12 to 23,
and weighted by the line.

Two evaluation samples are reported, because they estimate different things. The \emph{policy
sample} (Table~\ref{tab:misspec}) is the cycles each rule's own policy lent, and its Brier score is
the policy score: the accuracy of the predictions behind the loans that rule made. It mixes
prediction with selection. The rule's lines decide who borrows, since a borrower whose available
credit is below one USDC sits the month out, and a default ends the borrower's record under the rule
that observed it, so the three rules score different cycles (\CyclesWeightedAssumed,
\CyclesUnweightedAssumed\ and \CyclesTierOnlyAssumed\ in the assumed world). The \emph{common panel}
(Table~\ref{tab:misspecpanel}) removes the selection: it is every (seed, borrower, month) in which
the borrower would seek a loan, scored on the outcome that loan would have had, both drawn by the
world before any rule acts (the simulation uses common random numbers), so the panel is the same
\PanelCyclesAssumed\ cycles for every rule and does not depend on the lending decision. Each rule's
prediction is its posterior for that borrower at that month, whether or not it lent: this scores the
rule out of its own sample, on the borrowers it declined and on the months after a default it
observed, where its posterior carries that default as evidence. Uncertainty is bootstrapped on one
footing: the clusters are borrower paths, identified by (seed, borrower) and pooled across seeds
(\MisspecClusters\ clusters), each world draws \MisspecBoot\ resamples of them with replacement, and
the same resample serves every rule and both samples, so the intervals on the Brier scores, on the
skill and on the paired differences between rules are comparable. The interval covers borrower
sampling only; the state the rules share within a seed (the budget, the epochs, the farm) is not
resampled.

Three findings. First, at a two-year horizon the repayment evidence improves on the prior alone by
a measurable but small amount, and the weights do not improve on the unweighted rule. On the common
panel in the assumed world the weighted rule's Brier score minus the prior's is
$\PanelDeltaTierOnlyAssumed\times10^{-6}$ (interval $[\PanelDeltaLoTierOnlyAssumed,
\PanelDeltaHiTierOnlyAssumed]\times10^{-6}$, about one part in \PanelGainTierOnlyAssumed\ of the
score), and the interval excludes zero in \PanelBetterThanTierOnly\ of the \MisspecWorlds\ worlds; its skill
against the base rate is \PanelSkillWeightedAssumed\% (interval
$[\PanelSkillLoWeightedAssumed,\PanelSkillHiWeightedAssumed]$\%) against \PanelSkillTierOnlyAssumed\%
for the prior alone, whose interval includes zero. Against the unweighted rule the panel difference
is $\PanelDeltaUnweightedAssumed\times10^{-6}$ (interval $[\PanelDeltaLoUnweightedAssumed,
\PanelDeltaHiUnweightedAssumed]\times10^{-6}$), its interval includes zero in \PanelTiedWithUnweighted\
of the \MisspecWorlds\ worlds, and the weighted rule scores better in \PanelWeightedBeatsUnweighted\
of them. On the matched policy sample, the cycles both rules lent, the differences in the assumed
world are $\MatchedDeltaUnweightedAssumed\times10^{-6}$ against the unweighted rule and
$\MatchedDeltaTierOnlyAssumed\times10^{-6}$ against the prior alone, both intervals including zero
(\MatchedCyclesUnweightedAssumed\ and \MatchedCyclesTierOnlyAssumed\ matched cycles). The gain is
small because a live account shows about fifteen trials
against a prior worth fifty to ninety and defaults end the record. Second, misspecification degrades
every rule alike: with the riskier population the realised-to-predicted ratio is
\RatioWeightedRiskier\ under the weights and \RatioTierOnlyRiskier\ under the prior alone, and after
the regime shift \RatioWeightedShift\ and \RatioTierOnlyShift; the evidence corrects a wrong prior
only as fast as trials accumulate, and the ratio in months 12 to 23 under the shift
(\RatioLateWeightedShift) shows how slowly. A wrong lateness ratio moves calibration by less
(\RatioWeightedLateHalf\ and \RatioWeightedLateDouble). Third, the weights' contribution is
elsewhere: the bought identity that runs the recycled-seed farm holds a line of \FarmLineUnweighted\
under the unweighted rule, the tier's cap, against \FarmLineWeighted\ under the weights, which is the
line of a fresh identity with no history (Table~\ref{tab:misspeclines}); the unweighted rule also
grants honest survivors larger lines (\MeanLineUnweightedAssumed\ against \MeanLineWeightedAssumed\
USDC at the last report), because it counts trials that risked nothing. The narrow claim of
Section~\ref{sec:evidence} stands: selected zero-cost histories receive zero weight. The broader
claims do not follow from this simulation: that manufactured history is worthless in general, or
that risk-weighted repayments improve prediction. At this horizon the policy's accuracy is nearly
its prior's accuracy, and an issuer should expect to be as wrong as its prior is when the population
or the hazard moves. What the panel does not validate: engineered histories with a small positive
cost, proper scores that weight by exposure, and the outcome-dependent weights of
Section~\ref{sec:evidence}, which the panel scores but does not correct for.

\begin{table}[ht]
  \centering\small
  \caption{The policy sample: three scoring rules in five worlds (\MisspecHonest\ honest borrowers,
  \MisspecSeeds\ seeds pooled), each scored on the cycles its own policy lent. Number of cycles;
  Brier score of the predicted cycle default probability against the realised default ($\times100$)
  with a 95\% cluster bootstrap interval over the \MisspecClusters\ borrower paths; skill against the
  realised base rate; ratio of realised to predicted defaults over all cycles, over months 12 to 23,
  and weighted by the line.}
  \label{tab:misspec}
  \footnotesize\setlength{\tabcolsep}{3.5pt}
  \begin{tabular}{llrrrrrr}
    \toprule
    World & Scoring & Cycles & Brier $\times100$ [95\%] & Skill & Ratio & Months 12--23 & Line-wtd \\
    \midrule
    \inputrows{data/table_misspecification.tex}
    \bottomrule
  \end{tabular}
\end{table}

\begin{table}[ht]
  \centering\small
  \caption{The common panel: the same rules scored on every cycle a borrower would have sought
  (\PanelCyclesAssumed\ per world, the same for every rule), on the outcome that loan would have had.
  Brier score ($\times100$) with its 95\% cluster bootstrap interval; skill against the base rate;
  and the paired difference of the weighted rule's Brier score from this rule's on the same
  resamples ($\times10^{6}$; negative favours the weighted rule).}
  \label{tab:misspecpanel}
  \footnotesize\setlength{\tabcolsep}{4pt}
  \begin{tabular}{llrrr}
    \toprule
    World & Scoring & Brier $\times100$ [95\%] & Skill & Weighted $-$ this, $\times10^{6}$ [95\%] \\
    \midrule
    \inputrows{data/table_misspecification_panel.tex}
    \bottomrule
  \end{tabular}
\end{table}

\begin{table}[ht]
  \centering\small
  \caption{The same runs: honest loss per exposure-year, mean line of surviving honest borrowers at
  the last report (USDC), and the line held at the end by a bought basic identity that runs the
  recycled-seed farm (USDC; the fresh-identity line is 6.10 and the tier's cap 15).}
  \label{tab:misspeclines}
  \begin{tabular}{llrrr}
    \toprule
    World & Scoring & Loss rate & Mean line & Farm line \\
    \midrule
    \inputrows{data/table_misspecification_lines.tex}
    \bottomrule
  \end{tabular}
\end{table}

\paragraph{Attackers} Identity buyers' best strategy is to build history for 22 or 23 months and
then draw everything. Table~\ref{tab:attack} and Figure~\ref{fig:attack} compare two worlds that differ
only in the true price of an identity. When caps respect Proposition~\ref{prop:ic}, profit is negative
and falls linearly with the number of identities. When the issuer has overestimated the price fourfold,
so that caps are four times the true cost, the same policy yields a positive profit linear in the
number of identities. The loss on issued lines stays below $n\,\mathrm{cap}_t$ in every case.

\begin{table}[ht]
  \centering\small
  \caption{Identity buyers by tier: line at the bust-out month and profit per identity at the true
  cost and at a quarter of it (USDC), and with 64 identities the loss on issued lines against the
  bound $64\,\mathrm{cap}_t$.}
  \label{tab:attack}
  \begin{tabular}{lrrrrrrr}
    \toprule
    Tier & Cost & Cap & Line at bust & Profit, cost $k_t$ & Profit, cost $k_t/4$ & Loss, 64 & Bound \\
    \midrule
    \inputrows{data/table_attack.tex}
    \bottomrule
  \end{tabular}
\end{table}

\begin{figure}[ht]
  \centering
  % Figure: attacker profit against identities bought (data/attacker_curves.csv).
\begin{tikzpicture}
  \begin{axis}[
      width=0.72\linewidth, height=5.6cm,
      xmode=log, log basis x=2, xmin=1, xmax=64,
      xtick={1,2,4,8,16,32,64}, xticklabels={1,2,4,8,16,32,64},
      xlabel={identities bought}, ylabel={attacker profit (USDC)},
      tick label style={font=\footnotesize}, label style={font=\small},
      legend style={font=\scriptsize, draw=none, at={(0.03,0.03)}, anchor=south west},
      legend cell align=left,
      scaled y ticks=false, yticklabel style={/pgf/number format/fixed, /pgf/number format/1000 sep={,}},
    ]
    \addplot[thick, green!45!black, mark=*, mark size=1.2pt] table[x=identities, y=basic_ic, col sep=comma] {data/attacker_curves.csv};
    \addlegendentry{basic, true cost}
    \addplot[thick, blue!65!black, mark=square*, mark size=1.2pt] table[x=identities, y=standard_ic, col sep=comma] {data/attacker_curves.csv};
    \addlegendentry{standard, true cost}
    \addplot[thick, red!70!black, mark=triangle*, mark size=1.4pt] table[x=identities, y=enhanced_ic, col sep=comma] {data/attacker_curves.csv};
    \addlegendentry{enhanced, true cost}
    \addplot[thick, dashed, green!45!black, mark=*, mark size=1.2pt] table[x=identities, y=basic_viol, col sep=comma] {data/attacker_curves.csv};
    \addlegendentry{basic, cost a quarter}
    \addplot[thick, dashed, blue!65!black, mark=square*, mark size=1.2pt] table[x=identities, y=standard_viol, col sep=comma] {data/attacker_curves.csv};
    \addlegendentry{standard, cost a quarter}
    \addplot[thick, dashed, red!70!black, mark=triangle*, mark size=1.4pt] table[x=identities, y=enhanced_viol, col sep=comma] {data/attacker_curves.csv};
    \addlegendentry{enhanced, cost a quarter}
    \addplot[dotted, gray, domain=1:64] {0};
  \end{axis}
\end{tikzpicture}

  \caption{Attacker profit against the number of identities bought, at the true identity cost (solid)
  and when identities cost a quarter of what the issuer assumed (dashed).}
  \label{fig:attack}
\end{figure}

% ---------------------------------------------------------------------------------------------
\section{Limitations}

\begin{itemize}
\item \emph{Identity cost is an input.} Proposition~\ref{prop:ic} is conditional on $k_t$. If
  identities become cheaper, or one identity can be used at many lenders, caps must fall with them;
  Figure~\ref{fig:attack} shows what happens otherwise.
\item \emph{Thin data.} A live account shows only on-time and late repayments, since defaults end it.
  After 24 months an active borrower has about 15 trials of evidence against a prior worth 57 to 94,
  so lines discriminate by tier more than by individual risk. A real issuer adds off-chain data, which
  this policy admits only through the tier and its prior.
\item \emph{The prior and the late ratio carry the accuracy.} The calibration of Section~\ref{sec:sim}
  is measured in the population the prior is fitted to, with lateness at the assumed ratio; under a
  riskier population or a regime shift every scoring rule, the weighted one included, is wrong by
  the same factor, and at a two-year horizon the evidence adds little predictive skill over the
  prior and the weights none over an unweighted rule (Tables~\ref{tab:misspec}
  and~\ref{tab:misspecpanel}). Fractional likelihood weights define a power posterior, not a
  calibrated posterior for the observation process; survivorship (defaults end the record),
  selective lending and outcome-dependent weights are not corrected for.
\item \emph{Exponential demand and full drawdown at default} are modelling choices that give a concave
  profit and a closed-form line.
\item \emph{Not modelled:} staleness of scores, the pool's liquidity limits and provisioning, and the
  oracle infrastructure that carries reports.
\end{itemize}

\section{Conclusion}

An issuer of credit lines to pseudonymous borrowers can be made safe for lenders by the contract and
useful to borrowers by its policy. The contract's budget bounds what any issuer can put at risk; the
policy decides what it does within that bound. Caps no larger than the cost of an identity bound the
profit from identity fraud by its cost, evidence weighted by risk borne closes the routes to evidence
that cost nothing, and a greedy allocation over a laminar matroid spends the budget optimally and
survives front-running. What the policy does not do, at a two-year horizon, is predict individual
default better than its prior; its accuracy is the prior's, and the binding inputs are the price of
an identity and the base rates of the population, which the issuer must measure rather than assume.

\paragraph{Disclosure of AI authorship} This paper was written by Claude Code, an AI agent made by
Anthropic, working with the project's human operator, who set the research direction. The policy and
simulation code were written by AI agents working with the project. All of it is public and
reproducible.

{\sloppy\paragraph{Data and code availability} The policy, its tests and the simulation, at commit
b725a85, are in \texttt{analysis/issuer\_policy} of
\url{https://github.com/scottonchain/microcredit-contract}; the script that extracts every table and
figure is in \url{https://github.com/scottonchain/microcredit-theory}.\par}

\bibliographystyle{elsarticle-harv}
\bibliography{references}

\end{document}
